% ===== tr-02 块1：第 3 章 pp.1--4（原文页 1--4）=====
% [原文 p.1]
\section*{摘要}

我们追溯自 1996 年 Strominger--Yau--Zaslow 猜想提出以来的进展与思考。特别地，我们意在解释该猜想如何引出了由我本人与 Siebert 所发展的代数几何纲领，这一纲领的目标是通过研究 Calabi--Yau 流形的退化来解释镜像对称。最后，我们概述热带曲线如何在镜像对称的故事中出现。\footnote{2000 年美国数学会主题分类（2000 Mathematics Subject Classification）：14J32。本工作部分受 NSF 基金（资助号 0505325）资助。}

% [Tutor Note: 原文标题页有两条脚注（MSC 14J32 与 NSF 资助 0505325），译文合并为摘要末的一条脚注；同页页脚另有 arXiv 预印本的版权占位符一行，未译。]

% [Tutor Note: 原文引言标题 "Introduction." 未编号，故译为不编号节；本章第 1 节从原文 p.4 起，编号由 \section 自动给出（3.1、3.2……）。]

\section*{引言}

到 1996 年夏天为止，镜像对称领域已经取得了一系列引人注目的成功。在 Candelas、de la Ossa、Greene 与 Parkes [8] 关于五次三维流形（quintic threefold）瞬子预言的先驱性工作之后，Batyrev [2] 为环面簇中的 Calabi--Yau 超曲面给出了一个强有力的镜像对称构造，随后由 Batyrev 与 Borisov [3] 推广到环面簇中的完全交。Kontsevich [37] 引入了他的曲线稳定映射概念并研究了其模空间，为枚举几何的一次新的繁荣铺平了道路。最终，这以多种不同的形式促成了五次三维流形 Gromov--Witten 不变量的数学计算 [13, 43, 6, 12]。在此期间，弦论与代数几何两方面的许多研究者阐明了镜像对称的很大一部分结构。

另一方面，当时我主要感兴趣的是 Calabi--Yau 流形自身的几何。关于镜像对称的许多结果似乎主要依赖于周围环面簇的信息；特别地，Givental 的工作 [13] 以及 Lian、Liu 与 Yau [43]、Bertram [6] 和 Gathmann [12] 的后续工作，总是在映入 $\PP^4$ 而非五次三维流形的稳定映射的模空间上进行计算；而在 Batyrev--Borisov 型构造中，尽管周围环面簇之间有着清晰的组合关系，Calabi--Yau 流形本身之间却没有明显的几何关系。因此，正是由于缺乏几何理解，我倾向于回避对镜像对称的思考。

% [原文 p.2]
1996 年 6 月，随着 Strominger、Yau 与 Zaslow [56] 发表论文《镜像对称就是 $T$ 对偶》（``Mirror Symmetry is $T$-duality''），这一局面发生了戏剧性的改变。他们基于弦论中的近期想法提出了一个非凡的提议，首次为镜像对称给出了几何解释。

让我非常粗略地概述这里的物理论证。20 世纪 90 年代中期弦论的进展引入了 Dirichlet 膜（Dirichlet brane），即 D-膜的概念。它们是时空的（带某些附加数据的）子流形，应当充当开弦的边界条件，也就是说，我们允许开弦在传播时其端点被约束在一张 D-膜上。记住，按照弦论，时空看起来像 $\R^{1,3}\times X$，其中 $\R^{1,3}$ 是普通时空、$X$ 是 Calabi--Yau 三维流形，于是我们可以把一张 D-膜分解为 $\R^{1,3}$ 的一个子流形与 $X$ 上的一个子流形的乘积。结果表明------这里做了极大的简化------$X$ 上有两类特别值得关注的子流形：全纯 D-膜，即带有全纯线丛的全纯子流形；以及特殊拉格朗日 D-膜，即带有平坦 $U(1)$ 丛的特殊拉格朗日子流形：

\begin{definition}{3.0.1}
设 $X$ 是 $n$ 维 Calabi--Yau 流形，$\omega$ 是 $X$ 上某个 Ricci 平坦度量的 Kähler 形式，$\Omega$ 是一个无处为零的全纯 $n$-形式。若子流形 $M\subseteq X$ 是拉格朗日的，即 $\dim_{\R} M=\dim_{\C} X$ 且 $\omega|_M=0$，并且此外还有 $\operatorname{Im}\Omega|_M=0$，则称 $M$ 是特殊拉格朗日的。
\end{definition}

% [Tutor Note: 定义 0.1 的 Lagrangian 条件里特意写出 dim_R M = dim_C X（即 M 的实维数等于 X 的复维数），它等价于 M 是实 n 维的；这正是后文 McLean 定理给出 dim H^1(M,R) = n、从而 SYZ 纤维是 n-环面的维数前提。]

镜像对称在物理中的起源提示：若 $X$ 与 $\check X$ 是一对互为镜像的 Calabi--Yau 流形，则用 $X$ 紧化时空所得的弦论，应当与用 $\check X$ 紧化所得的弦论相同，只是其中某些数据被互换。就 D-膜而言，这一提示是：$X$ 上全纯 D-膜的模空间应当同构于 $\check X$ 上特殊拉格朗日 D-膜的模空间。而 $X$ 本身正是 $X$ 上各点的模空间。于是 $X$ 上的每个点应当对应于一个对 $(M,\nabla)$，其中 $M\subseteq\check X$ 是特殊拉格朗日子流形，$\nabla$ 是 $M$ 上的平坦联络。

McLean 的一个定理 [45] 告诉我们，特殊拉格朗日子流形 $M\subseteq\check X$ 的特殊拉格朗日形变模空间的切空间是 $H^1(M,\R)$。当然，$M$ 上平坦 $U(1)$-联络模去规范等价所得的模空间就是环面 $H^1(M,\R)/H^1(M,\Z)$。为使该模空间具有正确的维数，我们需要 $\dim H^1(M,\R)=n$，即 $X$ 的复维数。这提示 $X$ 由一族环面构成，而这族环面对偶于 $\check X$ 上的一族特殊拉格朗日环面。对这一论证略作细化，便得到如下猜想：

\begin{conjecture}{3.0.2}
（Strominger--Yau--Zaslow 猜想）若 $X$ 与 $\check X$ 是一对互为镜像的 Calabi--Yau $n$-流形，则存在纤维化 $f:X\to B$ 与 $\check f:\check X\to B$，其纤维是特殊拉格朗日的，且一般纤维是 $n$-环面。此外，这两个纤维化是对偶的，其意义是：只要 $X_b$ 与 $\check X_b$ 是非奇异环面，就典范地有 $X_b=H^1(\check X_b,\R/\Z)$ 与 $\check X_b=H^1(X_b,\R/\Z)$。
\end{conjecture}

% [Tutor Note: 原文此猜想句作 "then there exists fibrations"，为英语语法笔误（应为 exist），数学内容不受影响；译文按复数理解。另注意此处对偶写成等号 X_b = H^1(ˇX_b, R/Z) 而非同构符号，即把环面与其第一上同调群等同起来，这是原文的写法。]

在回顾过去十年的工作时，我会逐步澄清这一陈述；然而，就我这里的表述而言，这个猜想很可能是错的。另一方面，还有一些更弱的版本，它们大概是对的。更好的是，这些更弱的陈述大概已为现代技术所及（当然需要大量艰苦的工作）。尽管如此，在拓扑与辛的层面上，对上述猜想的一些精确版本已经取得了许多良好的进展。此外，该猜想在解释镜像对称的许多特征方面是成功的，其中有些特征仍只是启发性的，有些则已是严格的。我相信，对镜像对称的最终令人满意的理解将从 SYZ 猜想流出，即使结果并不采取它最初所提示的形式。

% [原文 p.3]
我在这里的主要目标是解释过去十年所走过的历程。我想把重点放在解释这些思想如何演变与发展，而不是放在精确的陈述上。除了最开始的几节之外，我只会给出很少的精确陈述。

在最初这几节里，我会澄清上面关于猜想的陈述，并说明它如何在所谓的半平坦情形（semi-flat case）下------即沿着特殊拉格朗日纤维的度量是平坦的情形------为镜像对称给出令人满意的解释。这自然引出关于仿射流形、其上的度量以及 Legendre 变换的讨论。这些如今看来正是支撑镜像对称的关键结构。

接下来我们考察出现奇异纤维的情形。此时我们需要放弃在半平坦情形中建立的那种精确形式的对偶性，而把注意力限制在拓扑对偶性上。在纯粹拓扑对偶性的范围内，SYZ 猜想在解释一大类 Calabi--Yau 流形（包括由 Batyrev--Borisov 构造给出的环面簇中的完全交）的镜像对称的拓扑特征方面是完全成功的。

再往下，我们考察 Dominic Joyce 的论证，它展示了上述强形式的 SYZ 猜想所面临的问题。这迫使我们将 SYZ 猜想改写为一个极限陈述。镜像对称总是与 Calabi--Yau 流形在极大幂零退化附近的行为有关。SYZ 猜想的一个极限形式提示：在极大幂零退化附近的 Calabi--Yau 流形上可以找到特殊拉格朗日环面；而当我们趋近复模空间中的极限点时，我们期待看到 Calabi--Yau 流形中越来越大的部分被特殊拉格朗日环面所填满。这一形式与原来的 SYZ 猜想不同，虽然仍然困难，但看上去有可能用现有技术触及。

% [Tutor Note: 原文 "semi-flat case, i.e. the case when the metric along the special Lagrangian fibres is flat" 中的 semi-flat 就是术语表里的「半平坦」；它指沿纤维方向度量（在纤维内）为平坦的拉格朗日纤维化，而非指整个 Calabi--Yau 度量为平坦。]

猜想的这一形式随后激发了一轮新的问题。在这一极限图景中，我们期待假设中的特殊拉格朗日纤维化的底 $B$ 是一列趋近极大幂零退化的 Calabi--Yau 流形的所谓 Gromov--Hausdorff 极限。Gromov--Hausdorff 收敛是一个度量空间概念，而极大幂零退化是一个代数几何的、Hodge 理论的概念。这两个概念如何关联？2000 年夏天，Kontsevich 提出：至少在拓扑层面上，该 Gromov--Hausdorff 极限粗略地说就是该代数几何退化的对偶相交复形。我们在第 6 节探讨这一想法。

另一方面，这对我们理解镜像对称有何帮助？与这些关于 SYZ 猜想极限形式的发展平行，我的合作者 Bernd Siebert 一直在与 Stefan Schröer 一起用对数几何研究 Calabi--Yau 流形的退化。Siebert 注意到，镜像对称似乎对应于一侧的对数数据与另一侧的极化之间的一个组合式的交换。我们一起意识到，这种处理镜像对称的途径与 SYZ 所预言的极限图景是吻合的。把这两种途径综合起来，我们发现了 SYZ 途径的一个代数几何版本，
% [原文 p.4]
下面我将加以描述。其基本想法是：忘掉特殊拉格朗日纤维化，只记住纤维化的底，那是一个仿射流形。我们说明仿射流形的多面体分解如何产生退化的 Calabi--Yau 簇；反之，某些类型的 Calabi--Yau 簇的退化（我们称之为环面退化）如何以其对偶相交复形的形式产生仿射流形。镜像对称再次由 Legendre 变换的一个离散版本加以解释，正如在半平坦情形中那样。

我们最后讨论热带曲线与这一途径的联系。用热带曲线在二维环面簇中计数曲线，是由 Mikhalkin [46] 的工作开创的；Nishinou 与 Siebert [48] 用一种直接受我这里所讨论的途径启发的办法，把这项工作推广到高维。另一方面，热带曲线迄今尚未被用于在 Calabi--Yau 流形中计数曲线，因此我将以讨论热带曲线如何在我们的图景中自然出现来结束本文。

% [Tutor Note: 原文 p.4 把 Mikhalkin 误拼作 "Mihkalkin"（500 dpi 裁图复核，p.4 该段第 2 行 "work of Mihkalkin [46]"），译文按通行拼写 Mikhalkin 处理；此处不涉及数学内容。]

我要感谢西雅图会议的组织者举办了一场出色的会议，也感谢我在 SYZ 相关问题上的合作者 Pelham Wilson 与 Bernd Siebert；本文提到的许多工作都出自与他们合作的工作。

% [Tutor Note: 本节（§1）起于原文 p.4。原文标题 "First a topological observation" 中的 observation 是「观察」而非「注记」，故译为「首先是一个拓扑观察」。]

\section{首先是一个拓扑观察}

在做任何别的事情之前，让我们先问一个非常基本的问题：为什么把环面纤维化取对偶会交换 Calabi--Yau 三维流形的 Hodge 数？如果你没见过这一点，它恰恰是最该首先考察的，因为只要作几条假设，它特别容易看出来。设给定一对 Calabi--Yau 三维流形 $X$ 与 $\check X$，带有纤维化 $f:X\to B$、$\check f:\check X\to B$，其性质是：存在稠密开集 $B_0\subseteq B$，使得 $f_0:f^{-1}(B_0)\to B_0$ 与 $\check f_0:\check f^{-1}(B_0)\to B_0$ 都是环面纤维丛。于是 $f$ 与 $\check f$ 的全部奇异纤维都位于 $\Gamma:=B\setminus B_0$ 之上。（注意：除非 $\chi(X)=0$，否则必定存在某些奇异纤维。）最后，假设 $f_0$ 与 $\check f_0$ 是对偶的环面纤维化，即 $f_0$ 可以等同于环面纤维化 $R^1\check f_{0*}(\R/\Z)\to B_0$，而 $\check f_0$ 可以等同于环面纤维化 $R^1f_{0*}(\R/\Z)\to B_0$。（这里记号略有滥用：所谓 $R^1\check f_{0*}(\R/\Z)$，我们真正指的是这样的环面丛------取与局部系统 $R^1\check f_{0*}\R$ 相伴的向量丛，再对格族 $R^1\check f_{0*}\Z$ 取商。）

% [Tutor Note: \check X 是原文用来表示镜像流形的记号（ˇX，check 重音），SYZ 文献惯例如此；不要误写为 \bar X 或 \tilde X。]

若 $V/\Lambda$ 是单个环面，$V$ 为 $n$ 维向量空间、$\Lambda$ 为 $V$ 中的格，则 $H^p(V/\Lambda,\R)\cong\bigwedge^pV^\vee$；而对偶环面 $V^\vee/\Lambda^\vee$ 有 $H^p(V^\vee/\Lambda^\vee,\R)\cong\bigwedge^pV$。若选取一个同构 $\bigwedge^nV\cong\R$，便得到同构 $H^p(V/\Lambda,\R)\cong H^{n-p}(V^\vee/\Lambda^\vee,\R)$。类似地，在 $f_0$ 与 $\check f_0$ 的相对情形下，若有一个同构 $R^3f_{0*}\R\cong\R$，就得到同构
\[
R^pf_{0*}\R\cong R^{3-p}\check f_{0*}\R.
\]

% [Tutor Note: 对偶环面方向上，H^p(V^∨/Λ^∨,R) ≅ ⋀^p V（注意这里是 V 而不是 V^∨），与 H^p(V/Λ,R) ≅ ⋀^p V^∨ 相对；两者维数都是 binom(n,p)，这是 Hodge 数交换的线性代数根源。]

现在我们作一条简化假设。设 $i:B_0\hookrightarrow B$ 为包含映射。若对所有的 $p$ 都有
\[
i_*R^pf_{0*}\R\cong R^pf_*\R,
\]
我们就说 $f$ 是 $\R$-simple 的。（一般地可以把 $\R$ 换成任意交换群 $G$，此时我们说 $f$ 是 $G$-simple 的。）当然，并非所有环面纤维化都是 $\R$-simple 的，但拓扑形式的 SYZ 中出现的最有意思的那些都是。于是设 $f$ 与 $\check f$ 是 $\R$-simple 的。在这一假设下，我们得到同构
\begin{equation}
R^pf_*\R\cong R^{3-p}\check f_*\R. \tag{3.1}
\end{equation}

% [Tutor Note: 式 (1)（译文本块唯一编号公式，编号 3.1）是 §1 的核心：它把 X 与 ˇX 的 Leray 谱序列的 E_2 项联系起来；注意指数 3-p 中的 3 是纤维化的底维数（这里假设 B 是三维流形），一般 n 维情形应为 n-p。]

% [Tutor Note: "R-simple" 中的 R 是实系数域的黑板体记号，不是斜体 R；G-simple 同理取任意交换群 G。原文 p.4 下文（下一块）用这一假设研究 f 与 ˇf 的 Leray 谱序列。]
